% Einheit 1 von 3 – Pyramide (bank/pyramide-kegel-kugel/e1.jsonl, zusammenbau v0.1)
%% TODO Zweigzeile Teil 1: was hier gelernt wird (Satz in Schülersprache aus dem Einheitstitel) – Einheit 1
\einheitenkopf[e1]{Einheit 1 von 3 · Pyramide}
\zweigzeile{ab Kl. 8, je nach Buch bis Kl. 9 · keine P10-Aufgabe}
\setcounter{aufgabe}{11}

%% TODO Titel: Ich-kann-Satz fehlt in der Bank (Platzhalter: Kettenname) – Nr. 12
%% TODO Anweisung: ein Satz über der ersten Teilaufgabe fehlt in der Bank – Nr. 12
\begin{aufgabe}{Säule, Spitze oder Kugel?}
\begin{teile}
\teil Zeltdach – Säule, Spitze oder Kugel? Kreuze an und färbe die Grundfläche.\\ \kreuz{Säule}\\ \kreuz{Spitze}\\ \kreuz{Kugel}

\pyramide{3.5}{3.5}{2.5}{}{}{}
\end{teile}
\end{aufgabe}

%% TODO Titel: Ich-kann-Satz fehlt in der Bank (Platzhalter: Kettenname) – Nr. 13
%% TODO Anweisung: ein Satz über der ersten Teilaufgabe fehlt in der Bank – Nr. 13
\begin{aufgabe}{Stützdreieck einzeichnen}
\begin{teile}
\teil Turmdach: Zeichne das rechtwinklige Dreieck aus Höhe, halber Grundkante und Seitenhöhe ein und markiere die Hypotenuse. Nichts rechnen.

\pyramide{4}{4}{3}{$6$ m}{}{$4$ m}
\end{teile}
\end{aufgabe}

%% TODO Titel: Ich-kann-Satz fehlt in der Bank (Platzhalter: Kettenname) – Nr. 14
%% TODO Anweisung: ein Satz über der ersten Teilaufgabe fehlt in der Bank – Nr. 14
\begin{aufgabe}{Volumen oder Oberfläche?}
\begin{teile}
\teil Stoff für ein pyramidenförmiges Zelt – Volumen oder Oberfläche? Kreuze an.\\ \kreuz{Volumen}\\ \kreuz{Oberfläche}
\end{teile}
\end{aufgabe}

%% TODO Titel: Ich-kann-Satz fehlt in der Bank (Platzhalter: Kettenname) – Nr. 15
%% TODO Anweisung: ein Satz über der ersten Teilaufgabe fehlt in der Bank – Nr. 15
\begin{aufgabe}{Welche Formel?}
\begin{teile}
\teil Pyramide mit der Grundfläche $30$ cm² und der Höhe $8$ cm, gesucht das Volumen – welche Formel? Schreibe sie aus der Formelsammlung ab und die Zahl zu jedem Buchstaben. Nichts rechnen.
\end{teile}
\end{aufgabe}

%% TODO Titel: Ich-kann-Satz fehlt in der Bank (Platzhalter: Kettenname) – Nr. 16
%% TODO Anweisung: ein Satz über der ersten Teilaufgabe fehlt in der Bank – Nr. 16
\begin{aufgabe}{Pyramide}
%% TODO \rechenplatz{n}: Schrittzahl des Grundfalls fehlt in der Bank (nur bei fester Schreibform, 2.3 b)
\begin{teile}
\teil Strecke von der Spitze senkrecht zur Mitte der Grundfläche – Höhe, Seitenhöhe oder Seitenkante? Kreuze an.\\ \kreuz{Höhe}\\ \kreuz{Seitenhöhe}\\ \kreuz{Seitenkante}
\teil Pyramide: Grundfläche $27$ cm², Höhe $8$ cm – Volumen? $V$ = \leerfeld[cm³]
\teil Quadratische Pyramide: Grundkante $6$ cm, Höhe $7$ cm – Volumen? $V$ = \leerfeld[cm³]
\teil Pyramide mit rechteckiger Grundfläche $8$ cm mal $5$ cm, Höhe $9$ cm – Volumen? $V$ = \leerfeld[cm³]
\teil Quadratische Pyramide: Grundkante $3{,}8$ cm, Höhe $5{,}5$ cm – Volumen? Runde auf eine Stelle. $V$ ≈ \leerfeld[cm³]
\teil Quadratische Pyramide: Grundkante $16$ cm, Höhe $15$ cm – Seitenhöhe? \leerfeld[cm]
\teil Quadratische Pyramide: Grundkante $8$ cm, Seitenhöhe $11$ cm – Flächeninhalt einer Seitenfläche? \leerfeld[cm²]
\teil Quadratische Pyramide: Grundkante $6$ cm, Seitenhöhe $9$ cm – Mantel? $M$ = \leerfeld[cm²]
\teil Quadratische Pyramide: Grundkante $8$ cm, Seitenhöhe $7$ cm – Oberfläche? $O$ = \leerfeld[cm²]
\teil Quadratische Pyramide: Grundkante $6$ cm, Höhe $7$ cm – Seitenkante? Runde auf zwei Stellen. \leerfeld[cm]
\teil Netz einer quadratischen Pyramide (verkleinert): Grundkante $4$ cm, Seitenhöhe $5$ cm. Zeichne in jede Seitenfläche die Seitenhöhe ein und schreibe die Maße an die Grundkanten und Seitenhöhen.

\netzpyramide{2}{2.5}{}{}
\teil Zeichne das Schrägbild einer quadratischen Pyramide mit Grundkante $4$ cm und Höhe $5$ cm (Tiefe halb so lang, schräg unter 45°). Zeichne die Höhe gestrichelt ein und beschrifte sie.

\rechenplatz{6}
\teil Pyramide: Volumen $96$ cm³, Grundfläche $36$ cm² – Höhe? $h$ = \leerfeld[cm]
\teil Ein Zeltdach hat die Form einer quadratischen Pyramide ohne Boden: Grundkante $5$ m, Seitenhöhe $3{,}2$ m. Wie viel m² Stoff braucht man? \leerfeld[m²]
\teil Das Dach eines Gartenhauses ist eine quadratische Pyramide mit der Grundkante $4{,}5$ m und der Höhe $1{,}8$ m. Zeichne in das Schrägbild die Höhe des Dachs ein. Beschrifte die Höhe und eine Grundkante mit den Maßen \hfill (P10 2015 OS).

\pyramide{4}{4}{2}{}{}{}
\end{teile}
\end{aufgabe}

%% TODO Titel: Ich-kann-Satz fehlt in der Bank (Platzhalter: Kettenname) – Nr. 17
%% TODO Anweisung: ein Satz über der ersten Teilaufgabe fehlt in der Bank – Nr. 17
\begin{aufgabe}{Pyramide – Fehler finden, Begründen, Darstellungswechsel, Anwendung}
\begin{teile}
\teil Pyramide: Grundfläche $48$ cm², Höhe $9$ cm. Jonas rechnet das Volumen: \rechnung{V &= 48 \cdot 9 \\ V &= 432 \text{ cm}^3} Wo ist der Fehler?
\teil Ein Quader und eine Pyramide haben dieselbe Grundfläche und dieselbe Höhe. Begründe, warum die Pyramide nur ein Drittel so viel Sand fasst.
\teil Aus diesem Netz wird eine Pyramide gefaltet. Skizziere ihr Schrägbild und beschrifte Grundkante und Seitenhöhe.

\netzpyramide{2}{3}{$4$ cm}{$6$ cm}
\teil Die Spitze eines Kirchturms ist eine quadratische Pyramide: Grundkante $6$ m, Seitenhöhe $11$ m. Die vier Dachflächen bekommen neues Kupferblech für $85$ € je m². Reichen $11\,000$ €?
\end{teile}
\end{aufgabe}
